ان (ت‌۲ + ت‌۴):}
  در فاصلهٔ ۱۳۵۸–۱۳۶۰، هر بحران بیرونی
  (بختیار، گروگان‌گیری، ترورها)
  فرصتی شد برای تمرکز بیشتر قدرت —
  و نخبگان انقلابی (بهشتی، رفسنجانی، خامنه‌ای)
  این مسیر را ساختارسازی کردند.

\item \textbf{تطور ایدئولوژیک نهادینه‌شده (ت‌۳ + ت‌۴):}
  از ۱۳۶۰ به بعد، مفاهیم «ولایت مطلقه» و «مصلحت نظام»
  جایگزین «هدایت معنوی» شدند
  و نظام سیاسی در قالبی نهادی تثبیت گشت
  که بازگشت از آن ناممکن بود.
\end{enumerate}

بنابراین، عنوان کتاب — «وعده یا خدعه؟» —
پرسشی دوگانه (\en{binary}) است
اما پاسخ \textbf{طیفی و مرحله‌ای} است.%
\footnote{%
\en{Keddie},
\textit{\en{Modern Iran}}, 2006, p.\,252:
``The truth almost certainly lies in a combination of factors.''}
\end{warningBox}

%============================================================
\section{جدول نهایی: پاسخ ترکیبی به تفکیک بزنگاه}
\label{sec:ch16-final-table}

جدول~\ref{tab:final-verdict} پاسخ ترکیبی را
به تفکیک هر بزنگاه خلاصه می‌کند.

{%
\small
\begin{longtable}{%
  >{\centering\arraybackslash}m{2.5cm}
  >{\centering\arraybackslash}m{2.8cm}
  >{\arraybackslash}m{7.5cm}}
\caption{پاسخ ترکیبی نهایی به تفکیک بزنگاه}
\label{tab:final-verdict} \\
\toprule
\textbf{بزنگاه}
  & \textbf{تز(های) غالب}
  & \textbf{توضیح مختصر} \\
\midrule
\endfirsthead
\caption[]{(ادامه)} \\
\toprule
\textbf{بزنگاه}
  & \textbf{تز(های) غالب}
  & \textbf{توضیح مختصر} \\
\midrule
\endhead
\midrule
\multicolumn{3}{r}{\footnotesize \textit{ادامه در صفحهٔ بعد}\,\dots} \\
\endfoot
\bottomrule
\endlastfoot

وعده‌های پاریس
  & \cellcolor{cAccent!10}ت‌۲ + خودفریبی
  & وعدهٔ عدم حاکمیت ناسازگار با \textit{ولایت فقیه} (۱۳۴۸)
    اما لحن مصاحبه‌ها حاکی از صداقت ذهنیِ نسبی است.%
    \footnote{صحیفهٔ امام، ج\,۵؛
    \en{Martin}, 2000, pp.\,78–95.} \\
\addlinespace

برخورد با بختیار
  & \cellcolor{cAccent!10}ت‌۲ + \cellcolor{cConsolid!10}ت‌۴
  & نصب بازرگان اما ایجاد ساختار موازی؛
    نخبگان انقلابی سازمان‌دهندهٔ اصلی.%
    \footnote{بازرگان، \textit{انقلاب ایران در دو حرکت}، ص\,۷۲.} \\
\addlinespace

مجلس خبرگان / قانون اساسی
  & \cellcolor{cAccent!10}ت‌۲ + \cellcolor{cGold!10}ت‌۳
  & بزنگاه نهادینه‌سازیِ ولایت فقیه؛
    هم فریب (وعده در برابر عمل)
    و هم تطور (گذار از هدایت به حاکمیت).%
    \footnote{\en{Schirazi}, 1997, pp.\,22–48.} \\
\addlinespace

ترورهای فرقان
  & \cellcolor{cConsolid!10}ت‌۴
  & بهره‌گیری نخبگان از بحران امنیتی
    برای تحکیم نهادهای موازی.%
    \footnote{\en{Abrahamian},
    \textit{\en{Radical Islam}}, 1989, pp.\,186–193.} \\
\addlinespace

بحران گروگان‌گیری
  & \cellcolor{cAccent!10}ت‌۲ + \cellcolor{cConsolid!10}ت‌۴
  & ابزارسازی آگاهانه از بحران؛
    تأیید ضمنی و سپس صریح؛
    رفراندوم قانون اساسی تحت سایهٔ بحران.%
    \footnote{\en{Bowden}, 2006, pp.\,204–219.} \\
\addlinespace

جنگ ایران و عراق
  & \cellcolor{cGold!10}ت‌۳ + \cellcolor{cConsolid!10}ت‌۴
  & تبدیل جنگ به ابزار بسیج و مشروعیت‌سازی؛
    تطور مفهوم «مصلحت نظام».%
    \footnote{\en{Crist}, 2012, pp.\,88–103.} \\
\addlinespace

حذف مجاهدین / بنی‌صدر / تصفیه‌ها / ۱۳۶۷
  & \cellcolor{cAccent!10}ت‌۲ + \cellcolor{cConsolid!10}ت‌۴
  & اوج تمرکز قدرت؛
    فتوای شخصی خمینی (ت‌۲)
    با اجرای نهادی نخبگان (ت‌۴).%
    \footnote{\en{Robertson}, 2010, pp.\,54–87;
    خاطرات منتظری، ۱۳۶۷.} \\

\end{longtable}
}

%============================================================
\section{نمودار نهایی: مسیر از وعده تا واقعیت}
\label{sec:ch16-final-tikz}

نمودار~\ref{fig:final-timeline} مسیر کلیِ حرکت خمینی
از وعده تا واقعیت را در یک خطِ زمانیِ واحد نشان می‌دهد
و تز غالب هر مرحله را مشخص می‌سازد.

\begin{figure}[ht]
\centering
\begin{tikzpicture}[
    event/.style={rectangle, rounded corners=4pt, draw=#1, thick,
                  fill=#1!8, minimum width=2.6cm, minimum height=1.5cm,
                  font=\tiny, align=center},
    yearnode/.style={circle, draw=cPrimary, fill=cPrimary!15,
                     font=\tiny\bfseries, inner sep=2pt,
                     minimum size=0.8cm},
    arr/.style={-{Stealth[length=5pt]}, thick, cGray},
  ]
  % main timeline
  \draw[very thick, cPrimary!40] (0,0) -- (15,0);

  % year markers
  \node[yearnode] (Y57) at (0,0) {۵۷};
  \node[yearnode] (Y58) at (3,0) {۵۸};
  \node[yearnode] (Y59) at (6,0) {۵۹};
  \node[yearnode] (Y60) at (9,0) {۶۰};
  \node[yearnode] (Y67) at (12,0) {۶۷};
  \node[yearnode] (Y68) at (15,0) {۶۸};

  % Phase 1
  \node[event=cAccent] (E1) at (1.5,2.5) {%
    \textbf{مرحلهٔ ۱}\\
    وعدهٔ پاریس\\[2pt]
    ت‌۲ + خودفریبی};
  \draw[arr, cAccent] (E1.south) -- ($(Y57)!0.5!(Y58)+(0,0.5)$);

  % Phase 2
  \node[event=cConsolid] (E2) at (6,2.5) {%
    \textbf{مرحلهٔ ۲}\\
    تمرکز قدرت\\[2pt]
    ت‌۲ + ت‌۴};
  \draw[arr, cConsolid] (E2.south) -- (Y59.north);

  % Phase 3
  \node[event=cGold] (E3) at (10.5,2.5) {%
    \textbf{مرحلهٔ ۳}\\
    نهادینه‌سازی\\[2pt]
    ت‌۳ + ت‌۴};
  \draw[arr, cGold] (E3.south) -- ($(Y60)!0.5!(Y67)+(0,0.5)$);

  % Power bar (bottom)
  \fill[cAccent!25] (0,-1.2) rectangle (15,-0.7);
  \shade[left color=cGreen!40, right color=cAccent!60]
    (0,-1.2) rectangle (15,-0.7);
  \node[font=\tiny, text=white] at (1.5,-0.95) {وعده};
  \node[font=\tiny, text=white] at (13,-0.95) {حاکمیت مطلقه};

  % bottom label
  \node[font=\scriptsize, text=cPrimary] at (7.5,-1.8)
    {شاخص تمرکز قدرت $\longrightarrow$};

  % Key events below timeline
  \node[font=\tiny, text=cGray, align=center] at (0,-2.5) {مصاحبه‌های\\پاریس};
  \node[font=\tiny, text=cGray, align=center] at (3,-2.5) {بختیار\\بازرگان};
  \node[font=\tiny, text=cGray, align=center] at (6,-2.5) {خبرگان\\گروگان‌گیری};
  \node[font=\tiny, text=cGray, align=center] at (9,-2.5) {حذف\\بنی‌صدر};
  \node[font=\tiny, text=cGray, align=center] at (12,-2.5) {کشتار\\۱۳۶۷};
  \node[font=\tiny, text=cGray, align=center] at (15,-2.5) {درگذشت\\خمینی};
\end{tikzpicture}
\caption{خط زمانی نهایی: از وعده تا حاکمیت مطلقه}
\label{fig:final-timeline}
\end{figure}

%============================================================
\section{پرسش‌های باز برای پژوهش آینده}
\label{sec:ch16-open}

هر پژوهش جدّی، پرسش‌هایی بی‌پاسخ بر جای می‌گذارد.
مهم‌ترین پرسش‌های باز عبارت‌اند از:

\begin{enumerate}
\item \textbf{اسناد آرشیوی:}
  آیا در آرشیوهای محرمانهٔ بیت رهبری یا دفتر رهبر انقلاب
  اسنادی وجود دارد که نشان‌دهندهٔ مباحثات داخلی
  دربارهٔ ولایت فقیه پیش از ۱۳۵۷ باشد؟%
  \footnote{میلانی، \textit{نگاهی به شاه}، مقدمه، ص\,۱۸
  بر ضرورت دسترسی آرشیوی تأکید دارد.}

\item \textbf{خاطرات منتشرنشده:}
  آیا یاران نزدیک خمینی
  (به‌ویژه احمد خمینی)
  خاطرات یا یادداشت‌هایی دارند که
  تردیدهای احتمالی خمینی را ثبت کرده باشند؟%
  \footnote{بخشی از بایگانی احمد خمینی
  در مؤسسهٔ تنظیم و نشر آثار امام نگهداری می‌شود
  اما به‌طور کامل در دسترس پژوهشگران نیست.}

\item \textbf{مقایسهٔ تطبیقی بین‌انقلابی:}
  آیا الگوی «وعدهٔ دموکراسی → تمرکز قدرت»
  در انقلاب‌های دیگر
  (روسیه ۱۹۱۷، کوبا ۱۹۵۹، نیکاراگوئه ۱۹۷۹)
  نیز تکرار شده است؟
  اگر بله، آیا مکانیزم‌های مشابهی در کارند؟%
  \footnote{%
  \en{Crane Brinton},
  \textit{\en{The Anatomy of Revolution}}, rev.\ ed., 1965, ch.\,8
  الگوی «ترمیدور» را برای انقلاب‌ها مطرح می‌کند.
  \en{Theda Skocpol},
  \textit{\en{States and Social Revolutions}}, 1979
  نیز تحلیل تطبیقی مشابهی ارائه می‌دهد.}

\item \textbf{نقش رسانه‌های غربی:}
  آیا رسانه‌های غربی با پذیرش بدون نقدِ وعده‌های پاریس
  عملاً در مشروعیت‌بخشی به خمینی نقش داشتند؟%
  \footnote{%
  \en{William Beeman},
  \textit{\en{The ``Great Satan'' vs.\ the ``Mad Mullahs''}},
  2005, ch.\,3.}

\item \textbf{مکانیزم‌های خودفریبی در رهبران سیاسی:}
  آیا یافته‌های عصب‌شناسیِ جدید دربارهٔ
  \en{self-deception}
  می‌تواند مدل دقیق‌تری برای رفتار خمینی ارائه دهد؟%
  \footnote{%
  \en{Sharot},
  \textit{\en{The Optimism Bias}}, 2011;
  \en{von Hippel \& Trivers},
  ``The Evolution and Psychology of Self-Deception'',
  \textit{\en{Behavioral and Brain Sciences}}, 34(1), 2011, pp.\,1–16.}

\item \textbf{پرسش هنجاری:}
  آیا در نظام حقوقیِ اسلامی (یا بین‌المللی)
  مکانیزمی برای پاسخگو کردن رهبران
  در قبال وعده‌های پیش‌انقلابی وجود دارد؟%
  \footnote{%
  \en{Thomas Franck},
  \textit{\en{The Power of Legitimacy Among Nations}}, 1990, ch.\,6.}
\end{enumerate}

%============================================================
\section{تأمل اخلاقی: مسئولیت وعده‌دهنده}
\label{sec:ch16-ethics}

فارغ از تبیین تاریخی، یک پرسش اخلاقی بنیادین باقی است:

\begin{warningBox}[پرسش اخلاقی]
آیا وعده‌دهنده‌ای که در بزنگاه‌های سرنوشت‌ساز
وعدهٔ عدم حاکمیت می‌دهد و سپس حاکم مطلق می‌شود،
از نظر اخلاقی مسئول است —
حتی اگر بخشی از تغییر مسیر ناخواسته بوده باشد؟
\end{warningBox}

سه سنّت اخلاقی پاسخ‌های متفاوتی می‌دهند:

\subsection{وظیفه‌گرایی کانتی (\en{Kantian Deontology})}

از منظر کانتی، \keyword{وعده‌شکنی} (\en{promise-breaking})
به‌خودی‌خود نادرست است،
صرف‌نظر از نتایج.%
\footnote{%
\en{Immanuel Kant},
\textit{\en{Groundwork of the Metaphysics of Morals}},
1785, 4:422;
ترجمهٔ فارسی: \textit{بنیاد مابعدالطبیعهٔ اخلاق}، ترجمهٔ عنایت‌الله شکیباپور.}
امرِ مطلق (\en{categorical imperative})
نمی‌پذیرد که وعده‌ای به‌عنوان ابزار تاکتیکی استفاده شود.
اگر خمینی آگاهانه وعده‌شکنی کرده باشد (ت‌۲)،
از نظر کانتی مسئول اخلاقی است.
اگر خودفریبی در کار بوده باشد،
مسئولیتش کمتر اما منتفی نیست —
زیرا \keyword{خودفریبی نیز صورتی از بی‌صداقتی}
با خود و دیگران است.%
\footnote{%
\en{Christine Korsgaard},
\textit{\en{Creating the Kingdom of Ends}}, 1996, pp.\,335–362.}

\subsection{پیامدگرایی (\en{Consequentialism})}

از منظر پیامدگرایانه، پرسش این نیست
که آیا وعده شکسته شد بلکه این است
که نتایج نهایی چه بودند.
اگر حاکمیت مطلقه به بهبود اوضاع جامعه انجامید،
وعده‌شکنی قابل توجیه خواهد بود.%
\footnote{%
\en{John Stuart Mill},
\textit{\en{Utilitarianism}}, 1863, ch.\,2;
ترجمهٔ فارسی: \textit{فایده‌گرایی}، ترجمهٔ مرتضی مردیها.}

اما شواهد تاریخی — از سرکوب ۱۳۶۰،
جنگ هشت‌ساله، کشتار ۱۳۶۷،
تا محدودیت آزادی‌های مدنی — نشان می‌دهد
که حتی از منظر پیامدگرایانه،
نتایج عمدتاً منفی بوده‌اند.%
\footnote{%
\en{Abrahamian},
\textit{\en{Tortured Confessions}}, 1999, pp.\,209–228;
\en{Robertson}, 2010, pp.\,1–15.}

\subsection{اخلاق فضیلت (\en{Virtue Ethics})}

از منظر \keyword{اخلاق فضیلت}
(\en{Aristotelian virtue ethics}),
پرسش این است: آیا خمینی فرد فضیلتمندی بود؟
فضیلت صداقت (\en{truthfulness})
و فضیلت عدالت (\en{justice})
هر دو اقتضا می‌کنند که حاکم به وعده‌هایش پای‌بند باشد —
نه فقط به‌دلیل وظیفه (کانت)
یا نتایج (پیامدگرایی)
بلکه به‌دلیل خصلت اخلاقی‌ای که حاکم باید داشته باشد.%
\footnote{%
\en{Alasdair MacIntyre},
\textit{\en{After Virtue}}, 3rd ed., 2007, ch.\,15.}

\begin{noteBox}[تبصرهٔ اسلامی]
در سنت اخلاقیِ اسلامی نیز
وفای به عهد (\textit{وفاء بالعهد})
از واجبات اخلاقی و فقهی شمرده می‌شود.%
\footnote{قرآن کریم، سورهٔ اسراء (۱۷)، آیهٔ ۳۴:
«وَأَوْفُوا بِالْعَهْدِ إِنَّ الْعَهْدَ كَانَ مَسْئُولاً.»}
مفهوم \keyword{مصلحت}
(\en{maṣlaḥa})
می‌تواند در شرایط استثنایی وعده‌شکنی را توجیه کند
اما گسترش بی‌حد آن خود مسأله‌ساز است —
چنان‌که تطور مفهوم مصلحت نظام در دههٔ ۱۳۶۰ نشان داد.%
\footnote{%
\en{Felicitas Opwis},
\textit{\en{Maṣlaḥa and the Purpose of the Law}}, 2010, pp.\,147–180.}
\end{noteBox}

%============================================================
\section{سخن پایانی}
\label{sec:ch16-final-word}

این کتاب تلاشی بود برای پاسخ به یکی از
مهم‌ترین پرسش‌های تاریخ معاصر ایران.
پاسخ نهایی — چنان‌که شاید قابل انتظار بود —
ساده نیست.
واقعیت تاریخی به‌ندرت در قالب دوگانه‌های مطلق می‌گنجد.

آنچه می‌توان با اطمینان نسبی گفت این است:

\begin{enumerate}
\item وعدهٔ خمینی در پاریس بدون تردید \textbf{نقض شد} —
  این واقعیتی تاریخی است، نه تفسیر.
\item \textbf{علت} نقض، ترکیبی مرحله‌ای از فریب،
  خودفریبی، فشار نخبگان و تطور ایدئولوژیک بوده است.
\item هیچ تز واحدی به‌تنهایی کفایت نمی‌کند
  و پژوهش آینده با اسناد جدید ممکن است
  وزن‌ها را تغییر دهد.
\item فارغ از تبیین تاریخی،
  \textbf{مسئولیت اخلاقی} وعده‌شکنی — حتی در ضعیف‌ترین خوانش —
  بر عهدهٔ وعده‌دهنده باقی می‌ماند.%
  \footnote{%
  این جمع‌بندی اخلاقی البته مسبوق به
  نقد منتظری بر خمینی است:
  «تاریخ شما را به‌عنوان یکی از [حذفی] ثبت خواهد کرد.»
  خاطرات منتظری، نامهٔ ۱۳۶۷.}
\end{enumerate}

%============================================================

\begin{principleBox}[خلاصهٔ نهایی کتاب — فصل ۱۶]
\begin{itemize}
\item \textbf{پاسخ ترکیبی مرحله‌ای:}
  \begin{enumerate}[nosep]
    \item مرحلهٔ اول (پاریس):
      \textcolor{cAccent}{فریب آگاهانه–نیمه‌آگاهانه}
      با عنصری از صداقت ذهنی.
    \item مرحلهٔ دوم (۱۳۵۸–۱۳۶۰):
      \textcolor{cAccent}{فریب} + \textcolor{cConsolid}{فشار نخبگان}.
    \item مرحلهٔ سوم (۱۳۶۰–۱۳۶۷):
      \textcolor{cGold}{تطور ایدئولوژیک} + \textcolor{cConsolid}{نخبگان}.
  \end{enumerate}

\item \textbf{وزن نهایی تزها (فصل~۱۵):}
  \textcolor{cAccent}{ت‌۲ ≈ ۳۳٪}،
  \textcolor{cConsolid}{ت‌۴ ≈ ۳۳٪}،
  \textcolor{cGold}{ت‌۳ ≈ ۲۳٪}،
  \textcolor{cGreen}{ت‌۱ ≈ ۱۱٪}.

\item \textbf{تز ۱ (صداقت)} ضعیف‌ترین اما نه کاملاً نادرست:
  مؤلفه‌ای جزئی در مرحلهٔ اول.

\item \textbf{الگوی قدرت:} یک‌سویه، انباشتی، بدون بازگشت.

\item \textbf{شش پرسش باز} برای پژوهش آینده شناسایی شدند
  (اسناد آرشیوی، خاطرات منتشرنشده، مقایسهٔ بین‌انقلابی،
  نقش رسانه‌ها، عصب‌شناسی خودفریبی، مکانیزم حقوقی پاسخگویی).

\item \textbf{تأمل اخلاقی:}
  از هر سه منظر وظیفه‌گرایی، پیامدگرایی و اخلاق فضیلت،
  وعده‌شکنی خمینی قابل توجیه کامل نیست.
  در سنت اخلاقی اسلامی نیز
  وفای به عهد واجب و مصلحت‌اندیشیِ بی‌حد مسأله‌ساز است.

\item \textbf{عنوان کتاب — «وعده یا خدعه؟»:}
  پرسش دوگانه است؛ پاسخ طیفی و مرحله‌ای.
  نه صرفاً وعده، نه صرفاً خدعه —
  بلکه فرایند پیچیده‌ای از فریب، خودفریبی،
  تطور و ساختارسازیِ نخبگانی
  که ایران را به مسیر کنونی‌اش رساند.
\end{itemize}
\end{principleBox}


%=============================================================
%  SEGMENT STATUS UPDATE
%=============================================================
%
%  Segment 0  (Preamble)           — COMPLETE
%  Segment 1  (Introduction, Ch 1-4) — COMPLETE
%  Segment 2  (Ch 5-8)             — COMPLETE
%  Segment 3  (Ch 9-12)            — COMPLETE
%  Segment 4  (Ch 13-16)           — COMPLETE  ✓
%
%  ALL 16 CHAPTERS + PREAMBLE PRODUCED.
%
%  Remaining optional tasks:
%    - Bibliography (BibLaTeX / .bib file)
%    - Index (makeidx / xindy)
%    - Appendices (source transcriptions, full timeline)
%    - Final proofreading pass
%    - Cover design
%
%=============================================================